\begin{resumo}
Este trabalho propõe e avalia um sistema automatizado de apoio à detecção de melanoma em imagens dermoscópicas, cujas principais contribuições metodológicas são (i) um \textbf{framework de cinco eixos} para avaliação automatizada da consistência estrutural e lógica de relatórios diagnósticos gerados por LLMs, e (ii) uma análise de \textbf{calibração de incerteza} que revela que o modelo de melhor desempenho concentra sua hesitação diagnóstica justamente nos casos clinicamente mais ambíguos. A arquitetura integra máscaras de segmentação de referência (\textit{ground truth}) do conjunto PH2 com agentes inteligentes baseados em modelos de linguagem multimodal (LLMs), os quais aplicam três protocolos diagnósticos clínicos independentes (ABCD, Menzies e SPCL (\textit{Seven-Point Checklist})) por meio de uma abordagem multiagente implementada com a biblioteca LangGraph. O conjunto PH2, composto por 200 imagens dermoscópicas anotadas por especialistas (40 melanomas e 160 não-melanomas), foi utilizado como base experimental, com suas máscaras de segmentação adotadas como entrada padronizada ao sistema para isolar a avaliação diagnóstica da variabilidade de modelos de segmentação automática. Foram avaliados os modelos GPT-5.4-nano (OpenAI) e Gemini-3.1-flash-lite (Google), comparados ao \textit{baseline} supervisionado ResNet-50. Os resultados indicam que as LLMs avaliadas apresentam desempenho quantitativo inferior ao \textit{baseline} convencional (ResNet-50, treinada diretamente sobre o PH2 com ponderação de classes: acurácia = 0,90; precisão = 0,57; sensibilidade = 1,00; F1 = 0,73) na maior parte das métricas, incluindo a sensibilidade — a métrica mais relevante para a detecção de melanoma, na qual deixar de identificar um caso maligno é o erro mais custoso, resultado esperado para modelos multimodais generalistas sem ajuste fino especializado. Mas a análise de consistência de cinco eixos revela um padrão informativo entre as próprias LLMs: o Gemini-3.1-flash-lite supera o GPT-5.4-nano em todos os eixos avaliados, incluindo aritmética ABCD correta (100\% vs.\ 96\%), aderência ao limiar declarado (46\% vs.\ 6\%), consenso inter-protocolos (78\% vs.\ 46\%) e acurácia nos casos com diagnóstico definitivo (85,2\% vs.\ 22,2\%). De forma decisiva, o Gemini-3.1-flash-lite demonstra calibração de incerteza clinicamente coerente, concentrando seus diagnósticos ``suspeitos'' e ``indeterminados'' nos casos com maior prevalência real de melanoma, enquanto o GPT-5.4-nano emite diagnóstico não-definitivo em 82\% dos casos sem discriminar adequadamente o nível de risco. Esses achados posicionam o framework de cinco eixos e a análise de calibração de incerteza como a contribuição metodológica central deste trabalho, indo além da acurácia classificatória bruta rumo a um protocolo reprodutível de auditoria da confiabilidade do raciocínio diagnóstico de LLMs. A proposta demonstra-se promissora como ferramenta de triagem assistida por IA para detecção de melanoma, especialmente em contextos com acesso limitado a especialistas em dermatologia.\looseness=-1

\textbf{Palavras-chave: Detecção de melanoma, Calibração de incerteza, Consistência diagnóstica, Modelos de linguagem multimodal, Sistema multiagente, PH2}
\end{resumo}
